\documentclass[../../../include/open-logic-section]{subfiles}

\begin{document}

\olfileid[fr]{his}{set}{pathology}
\olsection{Pathologies}

La définition de la \emph{limite} permit toutefois des constructions
assez «~pathologiques~».

Vers les années 1830, Bolzano découvrit une fonction
\emph{continue partout}, mais \emph{dérivable nulle part}.
Malheureusement, il ne publia jamais cette découverte : les
mathématiciens ne rencontrèrent cette idée qu'en 1872, lorsque
Weierstrass la redécouvrit indépendamment.\footnote{L'histoire est
documentée par les notes de bas de page, très détaillées, de
l'article de Wikipédia sur
\href{http://en.wikipedia.org/wiki/Weierstrass_function}{la
fonction de Weierstrass}.} Le résultat était pour le moins
surprenant. Il est facile de trouver des fonctions, telles que
$|x|$, qui sont continues partout mais non dérivables en un point
particulier. Une fonction continue partout et dérivable
\emph{nulle part} est d'une tout autre nature. Essayons un instant
d'imaginer comment la tracer : pour qu'elle soit continue, il doit
être possible de la dessiner sans jamais lever le crayon ; mais,
pour qu'elle ne soit dérivable en aucun point, il faudrait en changer
sans cesse et brusquement la direction.

D'autres «~pathologies~» suivirent. Le 5 janvier 1874, Cantor posa
à Dedekind, dans une lettre, le problème suivant :
\begin{quote}
Peut-on établir une correspondance biunivoque entre une surface
(par exemple un carré avec son bord) et une ligne (par exemple un
segment avec ses extrémités), de telle sorte qu'à chaque point de
la surface corresponde un point de la ligne et, réciproquement,
qu'à chaque point de la ligne corresponde un point de la surface ?

Pour le moment, il me semble encore très difficile de répondre à
cette question, bien qu'ici aussi l'on soit si fortement porté à
répondre \emph{non} que l'on aimerait considérer la démonstration
comme presque superflue. [Cité dans \citealt{Gouvea2011}.]
\end{quote}
Mais, en 1877, Cantor démontra qu'il s'était trompé. Une ligne et
un carré ont en fait exactement le même nombre de points. Le
29 juin 1877, il écrivit à Dedekind :
«~\emph{je le vois, mais je ne le crois pas}~». Dans le
«~récit reçu~» de l'histoire des mathématiques, ces mots sont
souvent présentés comme le signe que ces nouveaux résultats
paraissaient \emph{littéralement incroyables} aux mathématiciens
de l'époque. La correspondance est présentée dans
\citet{Gouvea2011}, et nous y reviendrons dans la
\olref[his][set][mythology]{sec}. La démonstration de Cantor est
esquissée dans la \olref[his][set][cantorplane]{sec}.

Inspiré par le résultat de Cantor, Peano se demanda s'il était
possible d'appliquer une ligne \emph{de façon continue} sur un
plan.\footnote{Précision de l'édition française : la source anglaise
emploie ici «~smoothly~», qui peut suggérer «~de manière lisse~».
Il faut comprendre «~de façon continue~» : l'image d'une courbe
lisse ne peut remplir une région plane.} Ce serait une
\emph{courbe remplissant l'espace}. En \citeyear{Peano1890},
Peano construisit précisément une telle courbe. Le résultat est
véritablement contraire à l'intuition : Euclide avait défini une
ligne comme une «~longueur sans largeur~» (Livre~I, définition~2),
mais Peano montra qu'en repliant convenablement une ligne, on peut
en quelque sorte transformer sa longueur en largeur. En
\citeyear{Hilbert1891}, Hilbert décrivit une courbe remplissant
l'espace un peu plus intuitive, accompagnée de dessins. Elle se
construit par étapes ; en voici les six premières :
\begin{center}
\begin{tikzpicture}
\begin{scope}[xshift=20pt, yshift=20pt]
	\draw [oldiagcolorC, l-system={Hilbert curve, step=40pt, angle=90, axiom=L, order=1}] lindenmayer system;
\end{scope}
\begin{scope}[xshift=110pt, yshift=10pt]
	\draw [oldiagcolorC, l-system={Hilbert curve, step=20pt, angle=90, axiom=L, order=2}] lindenmayer system;
\end{scope}
\begin{scope}[xshift=205pt, yshift=5pt]
	\draw [oldiagcolorC, l-system={Hilbert curve, step=10pt, angle=90, axiom=L, order=3}] lindenmayer system;
\end{scope}
\begin{scope}[xshift=2.5pt, yshift=-97.5pt]
	\draw [oldiagcolorC, l-system={Hilbert curve, step=5pt, angle=90, axiom=L, order=4}] lindenmayer system;
\end{scope}
\begin{scope}[xshift=101.25pt, yshift=-98.75pt]
	\draw [oldiagcolorC, l-system={Hilbert curve, step=2.5pt, angle=90, axiom=L, order=5}] lindenmayer system;
\end{scope}
\begin{scope}[xshift=200.625pt, yshift=-99.375pt]
	\draw [oldiagcolorC, l-system={Hilbert curve, step=1.25pt, angle=90, axiom=L, order=6}] lindenmayer system;
\end{scope}
\draw[gray] (0pt,0pt) rectangle (80pt, 80pt);
\draw[gray] (100pt,0pt) rectangle (180pt, 80pt);
\draw[gray] (200pt,0pt) rectangle (280pt, 80pt);
\draw[gray] (0pt,-100pt) rectangle (80pt, -20pt);
\draw[gray] (100pt,-100pt) rectangle (180pt, -20pt);
\draw[gray] (200pt,-100pt) rectangle (280pt, -20pt);
\end{tikzpicture}
\end{center}
À la limite, notion désormais définie rigoureusement, le carré
entier apparaît uniformément !!{colorC}. Au passage, la courbe de
Hilbert est continue partout mais dérivable nulle part. Comme on le
verra dans la \olref[his][set][hilbertcurve]{sec}, elle remplit un carré
de côté~$2^{-n}$ pendant chaque intervalle de durée~$4^{-n}$ de
l'étape~$n$. Fixons un instant~$t$ et choisissons un de ces intervalles
qui le contient. Deux sommets voisins du carré sont distants
de~$2^{-n}$. L'un d'eux est donc à une distance d'au moins~$2^{-n-1}$
de~$h(t)$. Il est atteint à un instant~$s_n$ du même intervalle, avec
$0<|s_n-t|\leq 4^{-n}$. La norme du quotient des accroissements est
ainsi au moins~$2^{n-1}$. Ces quotients ne sont pas bornés quand
$s_n$ tend vers~$t$, ce qui exclut une dérivée en~$t$.

Pour le meilleur ou pour le pire, ces constructions géométriques
«~pathologiques~» furent perçues comme une raison de se méfier de
l'intuition géométrique. Elles devinrent presque une
\emph{propagande} en faveur d'une nouvelle manière de faire des
mathématiques, qui allait culminer dans la théorie des ensembles.
Dans le récit que la discipline bâtit ensuite sur sa propre histoire,
on répéta souvent que ces résultats étaient à la fois parfaitement
rigoureux et profondément déconcertants. Ils servaient donc un double
but : mettre en garde contre le recours à l'intuition géométrique
et montrer la fécondité de nouvelles façons de penser.

\end{document}
